\begin{proof}
By  Propositions \ref{PropNSFWeightCast} and \ref{PropWAstLift}, $\hat{\varphi}_u^\natural$ is a proper approximate KMS-weight. Therefore, we can apply \cite[Theorem 3.4.5]{Des}, so that we obtain a measure $\mu^{M_1}$ on $\IR(\hat{N}_u)$, a measurable f\/ield of Hilbert spaces $(\mathcal{K}_\sigma^{M_1})_{\sigma \in \IR(\hat{N}_u)}$, a measurable f\/ield of representations $(\pi_\sigma)_{\sigma \in \IR(\hat{N}_u)}$, a measurable f\/ield of self-adjoint, strictly positive operators $(D_\sigma^{M_1})_{\sigma \in \IR(\hat{N}_u)}$ and an isomorphism~$\mathcal{Q}_0^{M_1}$ of $P_\gamma P_\beta \cHil$ onto $\int_{\IR(\hat{N}_u)}^\oplus \mathcal{K}_\sigma^{M_1} \otimes \overline{\mathcal{K}_\sigma^{M_1}} d \mu^{M_1}(\sigma)$ satisfying the properties of this theorem.



Let $\rho \in \IR(\hat{N}_u)$ be in the support of $\mu^{M_1}$.
We claim that $\pi_\rho$ is weakly contained in $\hat{\pi}_u \iota_u$. Suppose that this is not the case, so that there exists $x \in \hat{N}_u$ such that $\pi_\rho(x) \not = 0$ but $\hat{\pi}_u \iota_u(x)  = 0$. Let $X = \{ \sigma \in \IR(\hat{N}_u) \mid \pi_\sigma(x) \not = 0 \}$. Then $X$ is an open neighbourhood of $\rho$. Moreover, it follows form \cite[Theorem 3.4.5]{Des} that $\mathcal{Q}_0^{M_1}$ intertwines $\hat{\pi}_u \iota_u$ with $\int_{\IR(\hat{N}_u)}^\oplus \pi_\sigma \otimes \overline{1_{\mathcal{K}_\sigma}} d\mu^{M_1}(\sigma)$ from which it follows that $\mu^{M_1}(X) = 0$. This contradicts the fact that $\rho$ is in the support of $\mu^{M_1}$, so $\pi_\rho$ is weakly contained in $\hat{\pi}_u \iota_u$.

Since $\hat{\pi}_u \iota_u = \hat{\pi}\vert_{\hat{N}_c} \hat{\vartheta}^\natural$, we see that $\rho$ is in $\IR(\hat{N}_c)$, where $\IR(\hat{N}_c)$ is considered as a subset of~$\IR(\hat{N}_u)$ by the inclusion (\ref{EqnRepSpaces}).
Now we use Theorem \ref{ThmCorrespondenceC} to identify $\IR(\hat{N}_c)$ with $\IR(\hat{M}_c, M_1)$, which we consider as a subspace of~$\ICMM$ by (\ref{EqnRepSpaces}).  We consider $\mu^{M_1}$ as a measure on~$\ICMM$ by def\/ining the complement of~$\IR(\hat{N}_c)$ in~$\ICMM$ to be negligible. Let $U_\sigma \in \ICMM$ denote the corepresentation corresponding to $\sigma \in \IR(\hat{N}_c)$. So,
$
\pi_\sigma = \pi_{U_\sigma}^\natural.
$
We write $D_{U_\sigma}^{M_1}$ for $D_\sigma^{M_1}$ and set $D_U = 0$ for $U \in \ICMM $ not in the support of $\mu^{M_1}$. We denote $\cHil_U$ for the corepresentation Hilbert space of $U \in \ICMM$, and we get $\cHil_{U_\sigma}^{M_1} = \mathcal{K}_\sigma^{M_1}$.  Therefore, since the support of $\mu$ is contained in $\IR(\hat{N}_c)$, we see that $\mathcal{Q}_0^{M_1}$ is a map from $P_\beta P_\gamma \cHil \rightarrow \int^\oplus_{\ICMM} \cHil_U^{M_1} \otimes \overline{\cHil_U^{M_1}} d\mu^{M_1}(U)$.



(\ref{ItemPGOne}). It follows from  Corollary \ref{CorRep}, Remark \ref{RmkInvSubs} and the properties of $D_\sigma^{M_1}$ described in (1) of \cite[Theorem 3.4.5]{Des}, that for $\omega \in \mathcal{I}_N$, the operator $(\tilde{\omega} \otimes \iota)(U) (D_U^{M_1})^{-1}$ is bounded and its closure is Hilbert-Schmidt for $\mu^{M_1}$-almost every $U \in \ICMM$.

(\ref{ItemPGTwoA}) and (\ref{ItemPGTwoB}). We make two observations. First note that by  Proposition \ref{PropNSFWeightCast}, for $\omega \in \mathcal{I}_N$,
$
\hat{\Lambda}_u \iota_u( \lambda_u^\natural(\omega))  = \xi(\tilde{\omega}).
$
Secondly, we have proved that for every $\rho$ in the support of $\mu^{M_1}$, there is a $U \in \ICMM$, such that $\pi_\rho = \pi_U^\natural$. Then (2) of \cite[Theorem 3.4.5]{Des} yields (\ref{ItemPGTwoA}) of the present theorem. The second observation also yields (\ref{ItemPGTwoB}). Note that by Proposition \ref{PropMeas}, we see that $(\cHil_U)_{U \in \ICMM}$ is a measurable f\/ield of Hilbert spaces of which $(\cHil_U^{M_1})_{U \in \ICMM}$ forms a~measurable f\/ield of subspaces. Here we used that $\hat{M}_u$ is separable.


To prove (\ref{ItemPGThree}), we make the following observations. First of all, by Remark \ref{RmkGNSCorrespondence}, we see that $\hat{\pi}_u \iota_u = \pi_{W_{\mathcal{E}}}^\natural$, where $W_{\mathcal{E}}$ denotes the restriction of the multiplicative unitary $W$ to the Hilbert space $\mathcal{E}$. Secondly,
\[
\int_{\ICMM}^\oplus  \pi_U^\natural \otimes 1_{ \overline{ \cHil_U^{M_1}}} d\mu^{M_1}(U) =
\pi_{\int^\oplus_{\ICMM} U \otimes 1_{\overline{\cHil_U^{M_1}}} d\mu(U)}^\natural,
\]
where we use Proposition \ref{PropMeas} to infer that the direct integral on the right hand side exists.
Hence, by Remark \ref{RmkAlsoUniv} we see that $W_\mathcal{E}$ and $\int^\oplus_{\ICMM} U \otimes 1_{\overline{\cHil_U^{M_1}}} d\mu^{M_1}(U)$ are equivalent. Moreover, we def\/ine $\mathcal{Q}^{M_1}$ to be the intertwiner as constructed in the proof of Proposition \ref{PropEqRepCorep}. Then, $\mathcal{Q}^{M_1}$~satisf\/ies~(\ref{ItemPGThree}).
